\documentclass[10pt,a4paper]{amsart}
\usepackage[margin=19mm]{geometry}
\usepackage{amsmath,amssymb,mathtools}
\usepackage[scr=rsfs,cal=euler]{mathalfa}
\usepackage{xfrac}
\usepackage[unicode,hidelinks,bookmarks=false]{hyperref}
\newcommand{\ie}{i.e.\ }
\newcommand{\cf}{cf.\ }
\newcommand{\resp}{respectively\ }
\newcommand{\Subsection}{\subsection}
\providecommand{\nobreakdash}{\nobreak\hbox{-}}
\setlength{\emergencystretch}{2em}
\multlinegap0pt
\usepackage{bm}
\usepackage{bbm}
\usepackage{dsfont}
\usepackage{xspace}
\usepackage{cancel}
\usepackage{mathtools}
\usepackage{amsthm}
\makeatletter
\@ifundefined{Theorem}{\newtheorem{Theorem}{Theorem}}{}
\@ifundefined{Lemma}{\newtheorem{Lemma}[Theorem]{Lemma}}{}
\makeatother
\title[$q$-Numerical Ranges and Spectral Sets]{$q$-Numerical Ranges and Spectral Sets}
\author{Ryan O'Loughlin, Jyoti Rani}
\date{Source: arXiv:2603.15536v1 (CC BY 4.0)}
\begin{document}
\maketitle

\noindent\emph{Source and licence.} $q$-Numerical Ranges and Spectral Sets. Version arXiv:2603.15536v1 (CC BY 4.0), distributed under the Creative Commons Attribution 4.0 International licence, as recorded in the arXiv licence metadata for this exact version and shown on the abstract page https://arxiv.org/abs/2603.15536v1. Full rights notice: licenses/arxiv-2603.15536-CC-BY-4.0.txt.\par\smallskip

\noindent\textit{Editorial scope.} Counted unit: the Theorem whose source label is \texttt{None} in the article (source0.tex lines 369--420, source label \texttt{None}), delivered together with the article's own complete original proof of it (source0.tex lines 422--630, one \texttt{proof} environment of 209 lines). No mathematics is written, repaired or re-derived: statement and proof are the article's own text line for line. Required context retained uncounted: gamma (lines 159--162); per (lines 260--264); maingamma (lines 311--317). Every definition, display and label of those lines is retained unchanged. Omitted from this package and not redistributed: 1--158, 163--259, 265--310, 318--368, 421--421, 631--678. Nothing inside a retained excerpt is omitted or altered: every retained body is delivered exactly as the article's lines are written, so no omission receipt of any kind accompanies this package. Citations: the retained excerpts carry no citation command at all, so no bibliography entry of the article is transcribed and no reference is redistributed. Reference adaptation: none was needed and none was made; every reference inside the delivered unit resolves inside it, so no unresolved reference or citation is produced. Printed slips kept literal: no printed word, symbol, hypothesis or number of the counted statement or of its proof was altered. Layout and class: the article's own class is \texttt{article} with packages including graphicx, amsmath, times, hyperref, multicol, amsthm, amssymb, amsfonts, latexsym, color, enumitem, halloweenmath, caption, subcaption; the wrapper is the lane's pinned \texttt{amsart} editorial wrapper at 10pt on A4, which restates the theorem-like environments the retained text needs and adds the article's own macro definitions that the retained text needs (none), with extra wrapper packages (bm, bbm, dsfont, xspace, cancel, mathtools). The delivered numbering is the unit's own. Assets: no retained span contains a figure environment, a graphics-inclusion command, a PDF-page inclusion, a table or a TikZ picture; no image file is copied into this package, the package declares no asset dependency and no image file is redistributed with it. Licence: version arXiv:2603.15536v1 of this article is distributed under the Creative Commons Attribution 4.0 International licence, as recorded in the arXiv licence metadata for this exact version and shown on the abstract page https://arxiv.org/abs/2603.15536v1. A full rights notice is delivered at licenses/arxiv-2603.15536-CC-BY-4.0.txt. Characters: every retained excerpt is pure ASCII, so the delivered engine, XeLaTeX with its default Latin Modern text font, reports no missing character.
\begin{equation}\label{gamma}
    \gamma(f) := - \int_0^L \lambda_{\min}( \mu ( \sigma (s), A ) ) 
f( \sigma(s) )\,ds , \quad |\partial \Omega| = L .
\end{equation}

\begin{Lemma}\label{per}
$$
\int_0^{2 \pi} h_{\Omega}( e ^{i \theta}) d \theta = |\partial \Omega|.
$$
\end{Lemma}

\begin{Theorem}\label{maingamma}
Let $W(A)$ have non-empty interior. Then for $\gamma$ as given in \eqref{gamma},
    $$
    \gamma(1) \leq - \frac{1}{\pi }\left(\frac{1}{w_{\Omega}(A) + 2w_{W(A)}}\right)^2 \min_{\sigma \in \partial \Omega}\rho (\sigma)   \left( |\partial \Omega| -  | \partial W(A)| \right),
    $$
    where $\rho (\sigma) $ is the radius of curvature at a point $ \sigma \in \partial \Omega$.
\end{Theorem}

\begin{Theorem}
For $0 < |q| < 1$, let $\Omega_q = \frac{W_q(A)}{q}$ have smooth boundary and let $W(A)$ have non-empty interior. For $\theta\in[0,2\pi]$ and $\varepsilon>0$, define
\[
M_{\theta,\varepsilon}
:=
\left\{
x\in H:\|x\|=1,\ 
\Re\langle e^{-i\theta}Ax,x\rangle
\ge h_{W(A)}(e^{i\theta})-\varepsilon
\right\},
\]
and set
\[
m(\theta)
:=
\lim_{\varepsilon\downarrow 0}
\sup_{x\in M_{\theta,\varepsilon}}
\sqrt{\|Ax\|^2-|\langle Ax,x\rangle|^2}.
\]
Then
\[
\gamma(1)\le
-C\min_{\sigma\in\partial\Omega_q}\rho(\sigma)
\frac{\sqrt{1-|q|^2}}{|q|}
\int_0^{2\pi} m(\theta)\,d\theta,
\]
where
\[
C=\frac1\pi\left(\frac{1}{w_{\Omega_q}(A)+2w_{W(A)}}\right)^2.
\]

















Furthermore setting $t \in (0, \infty)$ such that $|q| = \frac{1}{\sqrt{1+t^2}}$ and writing $ \Omega(t) =  \Omega_{q(t)}$ we have 
$$
\int_0^{2 \pi} m(\theta) = \frac{d}{dt}_{|_{t=0+}}\left|\partial\Omega(t) \right|.
$$
\end{Theorem}

\begin{proof}
First observe that $\Omega_q = \Omega_{|q|}$ so it suffices to prove the theorem for $\Omega_{|q|}$. Furthermore, in light of the previous theorem and Lemma \ref{per} it suffices to prove that 
$$
\int_0^{2 \pi} \left( h_{\Omega_{|q|}}(e^{ i \theta}) - h_{W(A)}(e^{ i \theta}) \right) ds   \geq \frac{\sqrt{1-|q|^2}}{|q|} \int_0^{2 \pi} m(\theta).
$$
Fix $\theta\in[0,2\pi]$ and $x \in \mathcal{H}$ with $\|x\|=1$. Then any $ y \in \mathcal{H}$ with $\|y\|=1$ and $\langle x,y\rangle=|q|$ can be written as
$y=|q|x+\sqrt{1-|q|^2}\,z$ with $z\perp x$ and $\|z\|=1$. Hence
\begin{align}
    &\Re\big(e^{-i\theta}\langle Ax,y\rangle\big)
= |q|\,\Re\langle e^{-i \theta}A x,x\rangle
+\sqrt{1-|q|^2}\,\Re\langle e^{-i \theta}A x,z\rangle \\
& \le |q|\,\Re\langle e^{-i \theta}A x,x\rangle + \sqrt{1-|q|^2}\,\|P_{x^\perp}(e^{-i \theta}A x)\|,
\end{align}
where $P_{x^\perp}$ is the projection onto the orthogonal complement of the span of $x$. Clearly equality is attained by choosing $z$ in the direction of $P_{x^\perp}(e^{-i \theta}A x)$ and so
\[
h_{W_{|q|}(A)}(e^{i\theta})
=\sup_{\|x\|=1}\Big(|q|\,\Re\langle e^{-i \theta}A x,x\rangle+\sqrt{1-|q|^2}\,\|P_{x^\perp}(e^{-i \theta}A x)\|\Big),
\]
which means
\begin{equation}\label{star1}
    h_{\Omega_{|q|}}(e^{i\theta})
=\sup_{\|x\|=1}\Big(\Re\langle e^{-i \theta}A x,x\rangle+\frac{\sqrt{1-|q|^2}}{|q|}\,\|P_{x^\perp}(e^{-i \theta}A x)\|\Big).
\end{equation}
For every $\varepsilon>0$, restricting the supremum to $x\in M_{\theta,\varepsilon}$ gives
\[
h_{\Omega_{|q|}}(e^{i\theta})
\ge
\sup_{x\in M_{\theta,\varepsilon}}
\left(
\Re\langle e^{-i\theta}Ax,x\rangle
+
\frac{\sqrt{1-|q|^2}}{|q|}
\|P_{x^\perp}(e^{-i\theta}Ax)\|
\right).
\]
Hence
\begin{align}
    &h_{\Omega_{|q|}}(e^{i\theta})-h_{W(A)}(e^{i\theta}) \\
&\ge
\sup_{x\in M_{\theta,\varepsilon}}
\left(
\Re\langle e^{-i\theta}Ax,x\rangle-h_{W(A)}(e^{i\theta})
+
\frac{\sqrt{1-|q|^2}}{|q|}
\|P_{x^\perp}(e^{-i\theta}Ax)\|
\right).
\end{align}
Since
\[
\Re\langle e^{-i\theta}Ax,x\rangle \ge h_{W(A)}(e^{i\theta})-\varepsilon
\qquad \text{for } x\in M_{\theta,\varepsilon},
\]
it follows that
\[
h_{\Omega_{|q|}}(e^{i\theta})-h_{W(A)}(e^{i\theta})
\ge
-\varepsilon+
\frac{\sqrt{1-|q|^2}}{|q|}
\sup_{x\in M_{\theta,\varepsilon}}
\|P_{x^\perp}(e^{-i\theta}Ax)\|.
\]
Using
\[
\|P_{x^\perp}(e^{-i\theta}Ax)\|^2
=
\|e^{-i\theta}Ax\|^2-|\langle e^{-i\theta}Ax,x\rangle|^2
=
\|Ax\|^2-|\langle Ax,x\rangle|^2,
\]
we obtain
\[
h_{\Omega_{|q|}}(e^{i\theta})-h_{W(A)}(e^{i\theta})
\ge
-\varepsilon+
\frac{\sqrt{1-|q|^2}}{|q|}
\sup_{x\in M_{\theta,\varepsilon}}
\sqrt{\|Ax\|^2-|\langle Ax,x\rangle|^2}.
\]
Letting $\varepsilon\downarrow 0$ yields
\[
h_{\Omega_{|q|}}(e^{i\theta})-h_{W(A)}(e^{i\theta})
\ge
\frac{\sqrt{1-|q|^2}}{|q|}\,m(\theta).
\]
Integrating over $[0,2\pi]$, it follows that
\[
\int_0^{2\pi}
\bigl(h_{\Omega_{|q|}}(e^{i\theta})-h_{W(A)}(e^{i\theta})\bigr)\,d\theta
\ge
\frac{\sqrt{1-|q|^2}}{|q|}
\int_0^{2\pi} m(\theta)\,d\theta.
\]
The desired bound for $\gamma(1)$ now follows from Theorem \ref{maingamma} and Lemma \ref{per}.











We now prove the furthermore statement. From \eqref{star1}, noting that $
t=\frac{\sqrt{1-|q|^2}}{|q|}$, 
we deduce
\begin{align}
&\frac{h_{\Omega(t)}(e^{i\theta})-h_{W(A)}(e^{i\theta})}{t} \\
&=
\sup_{\|x\|=1}
\left(
\frac{\Re\langle e^{-i\theta}Ax,x\rangle-h_{W(A)}(e^{i\theta})}{t}
+
\|P_{x^\perp}(e^{-i\theta}Ax)\|
\right).
\end{align}
Fix $\varepsilon>0$. If $x\notin M_{\theta,\varepsilon}$, then $
\Re\langle e^{-i\theta}Ax,x\rangle-h_{W(A)}(e^{i\theta})<-\varepsilon$, 
and hence
\[
\frac{\Re\langle e^{-i\theta}Ax,x\rangle-h_{W(A)}(e^{i\theta})}{t}
+
\|P_{x^\perp}(e^{-i\theta}Ax)\|
<
-\frac{\varepsilon}{t}+\|A\|.
\]
Thus, for sufficiently small $t$, the supremum may be restricted to vectors
$x\in M_{\theta,\varepsilon}$, and so
\[
\frac{h_{\Omega(t)}(e^{i\theta})-h_{W(A)}(e^{i\theta})}{t}
\le
\sup_{x\in M_{\theta,\varepsilon}}
\|P_{x^\perp}(e^{-i\theta}Ax)\|.
\]
Taking $\limsup$ as $t\downarrow0$ gives
\[
\limsup_{t\downarrow0}
\frac{h_{\Omega(t)}(e^{i\theta})-h_{W(A)}(e^{i\theta})}{t}
\le
\sup_{x\in M_{\theta,\varepsilon}}
\|P_{x^\perp}(e^{-i\theta}Ax)\|.
\]
Since this holds for every $\varepsilon>0$, it follows that
\begin{equation}\label{eq1}
    \limsup_{t\downarrow0}
\frac{h_{\Omega(t)}(e^{i\theta})-h_{W(A)}(e^{i\theta})}{t}
\le
m(\theta).
\end{equation}

Let $\delta>0$. If $x\in M_{\theta,\delta t}$, then $
\Re\langle e^{-i\theta}Ax,x\rangle-h_{W(A)}(e^{i\theta})\ge -\delta t$ and so
\[
\frac{h_{\Omega(t)}(e^{i\theta})-h_{W(A)}(e^{i\theta})}{t}
\ge
-\delta+\|P_{x^\perp}(e^{-i\theta}Ax)\|.
\]
Taking the supremum over $x\in M_{\theta,\delta t}$ yields
\[
\frac{h_{\Omega(t)}(e^{i\theta})-h_{W(A)}(e^{i\theta})}{t}
\ge
-\delta+
\sup_{x\in M_{\theta,\delta t}}
\|P_{x^\perp}(e^{-i\theta}Ax)\|.
\]
Taking $\liminf$ as $t\downarrow0$, we obtain
\[
\liminf_{t\downarrow0}
\frac{h_{\Omega(t)}(e^{i\theta})-h_{W(A)}(e^{i\theta})}{t}
\ge
-\delta+m(\theta).
\]
Since $\delta>0$ is arbitrary, it follows that
\begin{equation}\label{eq2}
    \liminf_{t\downarrow0}
\frac{h_{\Omega(t)}(e^{i\theta})-h_{W(A)}(e^{i\theta})}{t}
\ge
m(\theta).
\end{equation}
Combining \eqref{eq1} and \eqref{eq2} gives
\[
\frac{d}{dt}\bigg|_{t=0+} h_{\Omega(t)}(e^{i\theta})=m(\theta).
\]

Observe that
\[
\|P_{x^\perp}(e^{-i\theta}Ax)\|^2
=
\|e^{-i\theta}Ax\|^2-|\langle e^{-i\theta}Ax,x\rangle|^2
=
\|Ax\|^2-|\langle Ax,x\rangle|^2,
\]
so
\[
0\le
\frac{h_{\Omega(t)}(e^{i\theta})-h_{W(A)}(e^{i\theta})}{t}
\le
\sup_{\|x\|=1}\|P_{x^\perp}(e^{-i\theta}Ax)\|
\le \|A\|.
\]
Thus the dominated convergence theorem allows us to integrate the derivative identity,
and by Lemma \ref{per} we obtain
\[
\int_0^{2\pi} m(\theta)\,d\theta
=
\frac{d}{dt}\bigg|_{t=0+} |\partial\Omega(t)|.
\]
\end{proof}

\end{document}
